\documentclass[11pt,a4paper]{article}
\usepackage[utf8]{inputenc}
\usepackage[T1]{fontenc}
\usepackage[spanish]{babel}
\usepackage{amsmath}
\usepackage{amssymb}
\usepackage{geometry}
\geometry{margin=2.5cm}

\title{\textbf{Introducción a la Cosmología \\ Guía 4: Nucleosíntesis}}
\date{}

\begin{document}
\maketitle

\section*{Problema 1}
Si los neutrinos electrónicos interactúan con protones y neutrones a temperaturas por encima de los $0.8\,\text{MeV}$ (freeze-out de los neutrones), ¿por qué decimos que están desacoplados a temperaturas por debajo de los $1.5\,\text{MeV}$?

\section*{Problema 2}
Encuentre la relación entre temperatura y tiempo durante la era dominada por la radiación, antes y después de la aniquilación electrón-positrón, en función del número de especies de neutrinos.

\section*{Problema 3}
Teniendo en cuenta que el deuterio tiene spin 1, obtenga la siguiente expresión para su concentración en masa en equilibrio:
\begin{equation}
  X_D=5.67\times10^{-14}\,\eta_{10}\,T_{\text{MeV}}^{3/2}\,e^{B_D/T}\,X_p\,X_n,
\end{equation}
donde $\eta_{10}=10^{10}\times n_N/n_\gamma$, $n_N$ es la densidad numérica de nucleones, $B_D=2.23\,\text{MeV}$ es la energía de ligadura del deuterio y $X_p$ y $X_n$ son las concentraciones de protones y neutrones libres respectivamente. Verifique que $X_D$ se hace de orden 1 a temperaturas cercanas a los $0.06\,\text{MeV}$. Obtenga también la concentración en equilibrio del helio-4, cuyo spin es 0 y cuya energía de ligadura es $28.3\,\text{MeV}$, y verifique que se hace de orden 1 a temperaturas cercanas a los $0.3\,\text{MeV}$. Usando estos resultados, explique qué es el cuello de botella del deuterio.

\section*{Problema 4}
Discuta cómo cambiaría la producción de helio-4 durante la nucleosíntesis primordial si una de las siguientes cantidades fuera mayor de lo que es: (a) la densidad de bariones; (b) la constante de acoplamiento de las interacciones débiles; (c) la constante de Newton; (d) la temperatura actual de la radiación cósmica de fondo.

\section*{Problema 5}
La figura a continuación muestra la concentración de neutrones libres en el universo en función de la temperatura y para distintos números de especies de neutrinos (fuente: Peebles, P.J.E., \emph{Principles of physical cosmology}, Princeton University Press, 1993).
\begin{enumerate}
  \item[(a)] ¿Por qué esta concentración tiene cambios en la pendiente? ¿A qué procesos corresponden?
  \item[(b)] ¿Por qué afecta el número de familias de neutrinos a la concentración de neutrones libres?
  \item[(c)] ¿A qué temperatura comienza a crecer apreciablemente la concentración de deuterio? Comparar esta temperatura con la energía de ligadura.
\end{enumerate}

\section*{Problema 6}
Sabiendo que la concentración de freeze-out de los neutrones es $X_n^*\simeq 0.158+0.005(N_\nu-3)$, donde $N_\nu$ es el número de especies de neutrinos, que la vida media de un neutrón libre es $\tau_n\simeq 886\,\text{s}$, que $\eta_{10}\simeq6$ y que la nucleosíntesis primordial empieza cuando la concentración en masa del deuterio es $X_D\simeq10^{-2}$, obtenga la concentración en masa del helio-4 al final de la nucleosíntesis en función de $N_\nu$.

\end{document}
